\section{A hyperbolic convex-hull realization}\label{sec:hyperbolic}

We apply Luo and Wu's Theorem~1.1(a) \cite{LuoWu2024} to the circle
domain supplied by Theorem~\ref{thm:main}. Identify $\Chat$ with the
ideal sphere $\partial_\infty\mathbb H^3$. A compact subset of that
sphere is \emph{circle type} if its complement is a circle domain.
Write $C(Y)$ for the hyperbolic convex hull of such a set $Y$.
For a three-dimensional hull, the intrinsic path distance on its boundary
is the infimum of hyperbolic lengths of boundary curves joining the points.

\begin{corollary}[Complete hyperbolic surfaces of genus zero]
\label{cor:hyperbolic}
Let $S$ be a connected orientable surface without boundary, of genus
zero, and let $g$ be a complete Riemannian metric on $S$ with Gaussian
curvature identically $-1$. There is a compact set
$Y\subset\partial_\infty\mathbb H^3$, containing at least three points,
whose complement is nonempty and connected, such that every connected
component of $Y$ is a closed round disk or a singleton and $(S,d_g)$
is isometric to $\partial C(Y)$ with its intrinsic path metric.

When $C(Y)$ is two-dimensional, $\partial C(Y)$ denotes the metric
double of $C(Y)$ across its relative boundary in the containing
totally geodesic plane. Surfaces are understood to be Hausdorff and
second countable; their number of ends is unrestricted.
\end{corollary}

\begin{proof}
The orientation and metric make $S$ a genus-zero Riemann surface.
Its classical planar realization, recalled in \cite[Section~1.1]{LuoWu2024},
and Theorem~\ref{thm:main} give a conformal equivalence from $S$ to a
circle domain $V$. The transported metric is complete with curvature
$-1$, so it is the unique complete conformal hyperbolic metric of $V$.
In particular, $V$ is hyperbolic and its boundary has at least three
points. Luo and Wu's Theorem~1.1(a), with their two-dimensional
convention in Section~1.1, now gives the asserted intrinsic isometry.
The realizing set has at least three points, since the convex hull of
at most two ideal points cannot supply a hyperbolic surface.
\end{proof}

The realizing set $Y$ need not be the complement of the intermediate
circle domain $V$. The corollary asserts existence and intrinsic
isometry; it supplies neither uniqueness nor a prescribed extension
to the ideal boundary. Its curvature hypothesis is exactly $-1$.
The realization of arbitrary complete metrics of curvature at least
$-1$ with each end represented by a circle or point at infinity, posed in
\cite[Section~1.2, Conjecture~3]{LuoWu2024}, is a separate problem.
